% Ручной файл (не генерируется): навигация по вопросам колобка.
% Источник вопросов: таблица «Методы Оптимизаций 2026», лист «(колобок) Теоремы и
% определения» (выгрузка 2026-10-27).
\chapter{Вопросы к колобку}\label{app:bilety}

\noindent Для каждого вопроса~--- кликабельные ссылки на заметки конспекта: сначала
то, что нужно рассказать, затем \emph{см.\ также}~--- определения и леммы, на
которые опирается ответ.

\begingroup
\setlist{leftmargin=3em}\sloppy
\section*{Теоремы и леммы}\label{app:bilety:thm}
\addcontentsline{toc}{section}{Теоремы и леммы}
\begin{description}[font=\normalfont\bfseries]
\item[Т1.] Свойства выпуклых множеств.\\
  \cref{note:08CF7113-4709-4C4F-A0E7-4A2ED595BD7F}~--- Выпуклость пересечения выпуклых множеств;\\
  \cref{note:80AFB48E-12EF-4D6C-9303-E2E52FC83D26}~--- Выпуклость проекции;\\
  \cref{note:88ABE812-A1C1-44A1-BF92-D714C9AF6315}~--- Выпуклость суммы Минковского выпуклых множеств;\\
  \cref{note:42DFF6DC-CC9D-40AA-9D6F-218DCF613833}~--- Характеристика выпуклого множества через выпуклые комбинации.
\item[Т2.] Внутреннее описание выпуклой оболочки.\\
  \cref{note:3A5B0E29-AED4-409E-A823-956F3A3957F7}~--- Внутреннее описание выпуклой оболочки.\\
  \emph{см.\ также:} \cref{note:42DFF6DC-CC9D-40AA-9D6F-218DCF613833}~--- Характеристика выпуклого множества через выпуклые комбинации; \cref{note:8583F986-A6E8-4BBB-B2FB-BA0E9BBE58BF}~--- Выпуклая комбинация.
\item[Т3.] Выпуклость в терминах касательных для дифференцируемых функций (дифф.\ критерий выпуклости).\\
  \cref{note:18BD8497-CEBD-4962-8144-474B489606C3}~--- Дифференциальный критерий выпуклости.\\
  \emph{см.\ также:} \cref{note:2F475EE4-2630-45FD-8DD2-8EA9D48059E4}~--- Дифференциальный критерий выпуклости функции одной переменной.
\item[Т4.] Дифференциальный критерий выпуклости (в терминах вторых производных).\\
  \cref{note:86FCE16B-C82E-4416-851C-0E0AD0FDA382}~--- Критерий выпуклости в терминах вторых производных.\\
  \emph{см.\ также:} \cref{note:9CEE854D-ED37-4C08-82EB-E41B36C43EA6}~--- Критерий выпуклости функции одной переменной через вторую производную.
\item[Т5.] Дифференциальный критерий оптимальности точек выпуклых функций.\\
  \cref{note:035B8935-65A9-4BB4-BDA3-60A60FF278F0}~--- Наименьшее значение выпуклой функции достигается в точке с нулевым градиентом.
\item[Т6.] Свойства выпуклых функций и операции, сохраняющие выпуклость.\\
  \cref{note:6BC7D685-C429-45A3-B563-269840CA0664}~--- Преобразования, сохраняющие выпуклость;\\
  \cref{note:1CD538BF-6F40-472A-A670-D38F093E4DF1}~--- Выпуклость множества подуровня выпуклой функции;\\
  \cref{note:D053F5F9-AF9F-4517-A434-44B68A40039F}~--- Выпуклость линии уровня для инфимума функции.\\
  \emph{см.\ также:} \cref{note:53558BAA-700E-46BE-9EAB-6A2F65A5B4C0}~--- Неравенство Йенсена.
\item[Т7.] Сильная выпуклость в терминах касательных для дифференцируемых функций.\\
  \cref{note:1B5A3E25-56E8-4D5E-9CF6-E2731DA77604}~--- Дифференциальный критерий сильной выпуклости.\\
  \emph{см.\ также:} \cref{note:51CC880D-BF19-43E2-9599-FB31A9A7367F}~--- Единственность минимума сильно выпуклой функции.
\item[Т8.] Дифференциальный критерий сильной выпуклости (в терминах вторых производных).\\
  \cref{note:94DADA28-87E4-4E36-B289-8415BCE750FC}~--- Критерий сильной выпуклости в терминах вторых производных.
\item[Т9.] Свойства функций с липшицевым градиентом. Связь с максимальным по модулю собственным числом матрицы вторых производных.\\
  \cref{note:379E4CAD-A44B-4631-8DAF-F3DC708EAEDE}~--- Усиленное неравенство для функций с липшицевым градиентом;\\
  \cref{note:4B6D53BA-E599-4DEA-A557-550BD32042BA}~--- Связь спектральной нормы гессиана и липшицевости градиента;\\
  \cref{note:DAA0DAE4-4B48-4A2F-9459-54790B6B129F}~--- Оценка нормы гессиана сверху для функции с липшицевым градиентом.\\
  \emph{см.\ также:} \cref{note:3C6CF60E-0773-488E-A0B9-4D7510D456C0}~--- Связь сильной выпуклости и липшицевости градиента; \cref{note:81C8B68E-DAAD-4B86-BAEB-AB0E1C07442F}~--- Аппроксимация функции с липшицевым гессианом до 3 степени.
\item[Т10.] Оценка значений функций с липшицевым градиентом через разность градиентов.\\
  \cref{note:655064DF-D686-4726-83DC-6F3D67CC5B9C}~--- Оценка для выпуклых функций с липшицевым градиентом.
\item[Т11.] Точная оптимизация. Метод золотого сечения. Метод парабол.\\
  \cref{note:8EDDABAC-D2BE-4EC9-8297-DBAC52F70D28}~--- Точная одномерная оптимизация;\\
  \cref{note:CD52D818-AA25-471C-97C2-5DBE9787CE53}~--- Унимодальная функция;\\
  \cref{note:1D154EFA-62D8-4BC3-B9CD-CCDE31F166B0}~--- Монотонность унимодальной функции;\\
  \cref{note:96DED7D9-3918-45C5-9850-E05030941645}~--- Метод золотого сечения;\\
  \cref{note:4E950058-62DB-4EDA-A71D-98CACBBC41AA}~--- Метод парабол.\\
  \emph{см.\ также:} \cref{note:1BCFB956-2053-49A8-ADDF-0D85D53F0610}~--- Метод Брента; \cref{note:124BD9BC-084E-43FB-B617-AF8690B33EF8}~--- Дихотомия по производной.
\item[Т12.] Неточная оптимизация. Условие Армихо--Вульфа. Существование.\\
  \cref{note:A818191F-B1A6-492C-B38D-5369D169161C}~--- Неточная одномерная оптимизация;\\
  \cref{note:5330B23D-95DD-4D5D-B4A3-63EBADB1DB7D}~--- Условия Армихо-Вульфа;\\
  \cref{note:C143295F-C422-4FEE-95AD-3D6B620BC0E8}~--- Существование шага для Армихо-Вульфа.\\
  \emph{см.\ также:} \cref{note:15611E05-13FC-4FB0-AE00-A4BB857A29E6}~--- Условие Армихо; \cref{note:2084FD2F-E2AA-4081-B9D1-14CEA27EA243}~--- Условие Вульфа; \cref{note:03E395A1-0606-49C1-978A-7275296FE023}~--- Алгоритмы для подбора шага с условиями Армихо-Вульфа; \cref{note:F16921E1-67C5-4EDE-A55F-AD11D8257A69}~--- Алгоритм Backtracking для подбора шага Армихо-Вульфа.
\item[Т13.] Сходимость градиентного метода для случая $C^{1,1}_L$.\\
  \cref{note:C3B64D64-D35B-42AC-B683-6736227CBBBC}~--- Сходимость градиентного спуска для функций с липшицевым градиентом;\\
  \cref{note:FFEA1F1D-B70C-4B4A-9386-8344A2A3BBC4}~--- Оценка убывания за шаг градиентного спуска.
\item[Т14.] Сходимость градиентного метода для случая $C^{1,1}_L$ и $\mu$-сильной выпуклости.\\
  \cref{note:85D9D77D-13A2-42E0-BF93-AD085E0134B9}~--- Сходимость градиентного спуска для сильно выпуклых функций с липшицевым градиентом.\\
  \emph{см.\ также:} \cref{note:3C6CF60E-0773-488E-A0B9-4D7510D456C0}~--- Связь сильной выпуклости и липшицевости градиента.
\item[Т15.] Сходимость градиентного метода с произвольным направлением с шагом по условию Армихо--Вульфа.\\
  \cref{note:36F6E1BA-41DB-460D-A9C6-B980621727D4}~--- Сходимость градиентного спуска с шагом Армихо-Вульфа;\\
  \cref{note:93FC9CC4-79FA-4CEA-9B10-9DBFB2C07850}~--- Оценка на шаг Армихо-Вульфа для градиентного спуска.\\
  \emph{см.\ также:} \cref{note:DF0CC2BA-9234-47D3-8988-7A1DAFC9B114}~--- Градиентный спуск с шагом Армихо-Вульфа; \cref{note:E4AE9B25-BDB9-4A54-8CCB-09809F20B79D}~--- Косинус угла между градиентом и направлением спуска.
\item[Т16.] Локальная сходимость метода Ньютона.\\
  \cref{note:CDA139F9-5438-41E1-96B1-11A751BE40C5}~--- Локальная сходимость метода Ньютона.\\
  \emph{см.\ также:} \cref{note:C0ED5501-CFC6-42BA-A53A-1B346CBA545D}~--- Метод Ньютона; \cref{note:FDA5DA6F-DE0F-437E-A4F4-89757BA7F538}~--- Квадратичная сходимость; \cref{note:4CD0DF77-0553-4E99-901D-4D4C50C9900C}~--- Квадратичная сходимость сверхлинейна.
\item[Т17.] Глобальная сходимость метода Ньютона.\\
  \cref{note:CC41F7D2-EECB-46E4-817A-756E7CB9B7C0}~--- Глобальная сходимость метода Ньютона.\\
  \emph{см.\ также:} \cref{note:C0ED5501-CFC6-42BA-A53A-1B346CBA545D}~--- Метод Ньютона.
\item[Т18.] Конструкция оценивающих последовательностей и сходимость <<спусков>>, построенных по ним.\\
  \cref{note:866FE52A-109C-4912-B61C-F29BBD475F6D}~--- Оценивающая последовательность;\\
  \cref{note:D21EBF43-7558-45EE-8198-60958E391E03}~--- Оценка невязки через оценивающую последовательность;\\
  \cref{note:C67AAFD8-D294-480B-B58E-C284BFF76D50}~--- Лемма о конструкции оценивающей последовательности;\\
  \cref{note:4B497C50-C541-428D-94F6-B0F0DF464A50}~--- Оценивающая последовательность с квадратичной начальной функцией.
\item[Т19.] Метод Нестерова.\\
  \cref{note:BB5E3FC2-069C-44F0-92D0-B4DC13092CCA}~--- Метод Нестерова;\\
  \cref{note:C038105E-BD42-44A8-9542-2DBE9FFAAA3F}~--- Сходимость ускоренного метода Нестерова.\\
  \emph{см.\ также:} \cref{note:C6A43F85-13E0-46B3-96A8-BCB9F34ECD57}~--- Метод тяжёлого шарика.
\item[Т20.] Условие Денниса--Море.\\
  \cref{note:D65FB534-B7F1-4194-8E4B-D4ACAA89E6D2}~--- Критерий Денниса–Море.
\item[Т21.] О сверхлинейной сходимости квазиньютоновских методов.\\
  \cref{note:A9E5C7FA-46FE-460E-A1B9-5B7CD72DB5B4}~--- Сверхлинейная сходимость квазиньютоновских методов.\\
  \emph{см.\ также:} \cref{note:6BA18878-DE34-46A5-ABE2-EDA09CAE56A7}~--- Квазиньютоновские методы.
\item[Т22.] Теорема о сходимости SGD для выпуклых функций.\\
  \cref{note:58D1735C-5A70-4DA0-B522-B4595CB4494C}~--- Сходимость SGD для выпуклых функций.\\
  \emph{см.\ также:} \cref{note:A4769F74-2E6F-4CA3-89E4-7B01881D8109}~--- Стохастический градиентный спуск (SGD).
\end{description}

\section*{Определения}\label{app:bilety:def}
\addcontentsline{toc}{section}{Определения}
\begin{description}[font=\normalfont\bfseries]
\item[О1.] Выпуклые множества, выпуклая оболочка.\\
  \cref{note:316B5BFA-8B9F-4267-9086-D6593AAB5F9B}~--- Выпуклое множество;\\
  \cref{note:1AC657F9-2A3D-47D1-A71D-5B813D961260}~--- Выпуклая оболочка;\\
  \cref{note:8583F986-A6E8-4BBB-B2FB-BA0E9BBE58BF}~--- Выпуклая комбинация.
\item[О2.] Выпуклые функции.\\
  \cref{note:1CA14F2C-637F-41DE-9E97-72075D19E97F}~--- Выпуклая функция.
\item[О3.] Линии уровня.\\
  \cref{note:BF27A4AF-11BF-4A0D-85DF-5F7790BD58A9}~--- Линия уровня функции.
\item[О4.] Сильно выпуклые функции.\\
  \cref{note:22E546FC-186F-4A67-85AA-2C85DCEF9021}~--- Сильно выпуклая функция.
\item[О5.] Функции с липшицевым градиентом.\\
  \cref{note:7DD1C7AB-1A49-4675-9F21-5A16C717F670}~--- Функции с липшицевым градиентом.\\
  \emph{см.\ также:} \cref{note:7E287F4B-AC3E-4CF5-9D32-3BC2BC52D01A}~--- L-гладкие функции.
\item[О6.] Виды сходимостей: сублинейная, линейная и сверхлинейная.\\
  \cref{note:A74C35E2-DDF7-4115-99B7-2448CF7C6878}~--- Сублинейная сходимость;\\
  \cref{note:DF2405BF-8B31-47EB-8DEF-B1A88034717C}~--- Линейная сходимость;\\
  \cref{note:9FA20DA5-3A95-4683-9CB5-C6C9A903C86A}~--- Сверхлинейная сходимость.\\
  \emph{см.\ также:} \cref{note:C42D2C19-B35F-4C0E-AFEB-D880CBB72176}~--- Невязка; \cref{note:55E52506-D4CB-4D20-AB2E-3793825A406F}~--- Скорость сходимости; \cref{note:FDA5DA6F-DE0F-437E-A4F4-89757BA7F538}~--- Квадратичная сходимость; \cref{note:C0408924-F658-43DF-8794-0C9AA85E7580}~--- Тест корней.
\item[О7.] Оценивающие последовательности.\\
  \cref{note:866FE52A-109C-4912-B61C-F29BBD475F6D}~--- Оценивающая последовательность.
\item[О8.] Квазиньютоновские методы.\\
  \cref{note:6BA18878-DE34-46A5-ABE2-EDA09CAE56A7}~--- Квазиньютоновские методы.\\
  \emph{см.\ также:} \cref{note:E976DE49-0091-45E1-B621-A236F8B87859}~--- Условие секущей; \cref{note:01306845-D37F-480A-8DCE-EE9CA022DA1A}~--- BFGS; \cref{note:F7506DC4-37FB-4E88-8178-B0FC678233C0}~--- SR1.
\end{description}
\endgroup
